%% Our colors (for backgrounds and code listings):
\definecolor{light-gray}{gray}{0.9}
\definecolor{dark-gray}{gray}{0.4}
\definecolor{light-blue}{RGB}{64,64,255}
\definecolor{dark-blue}{RGB}{16,16,64}
\definecolor{dark-green}{RGB}{16,128,16}

%Snippet for the circuits
\newcommand{\GridCircuitZoneFault}{
    The grid model and its operational point is as in the following schematic:
    \begin{center}
        \includestandalone[width=0.8\textwidth]{circuit_z1_fault.tikz}
    \end{center}
    \begin{center}
        \small \textbf{Note: This schematic is only a reminder of the test setup on the TSO's
        side --- the Producer's side may vary, depending on the user-provided model.}
    \end{center}
}

\newcommand{\GridCircuitZoneSetPoint}{
    The grid model and its operational point is as in the following schematic:
    \begin{center}
        \includestandalone[width=0.8\textwidth]{circuit_z1_setpoint.tikz}
    \end{center}
    \begin{center}
        \small \textbf{Note: This schematic is only a reminder of the test setup on the TSO's
        side --- the Producer's side may vary, depending on the user-provided model.}
    \end{center}
}
%End snippet for the circuits

%Snippet for the simulation section header
% #1 = description items (one or more \GridCurvesXxxDescription commands)
\newcommand{\GridSimulationHeaderZone}[1]{
    \subsubsection{Simulation}
    As required by the DTR PCS \DTRPcs, the figures below show the
    following magnitudes:
    \begin{itemize}
        #1
    \end{itemize}

    \subsubsection{Simulation results}
    The blue line shows the calculated curve and the orange line shows the reference curve.
}
%End snippet for the simulation section header

%Snippet for the figures of a test and their descriptions. The report shows the figures the
% tool drew for the test, or left a placeholder for, in the order below and two per row, each
% with its item in the description list: which ones there are depends on the test.
% #1 = figure name, #2 = what goes in the report when that figure is there for the test
\newcommand{\IfFigureDrawnZone}[2]{%
    \IfFileExists{\Producer_#1_\PCSname.\DTRPcs.\OCname.pdf}{#2}{}%
}

% #1 = figure name, #2 = description item
\newcommand{\GridCurveDescriptionZone}[2]{\IfFigureDrawnZone{#1}{\item #2}}

\newcounter{gridfigure}

% #1 = figure name, #2 = caption. The figure takes the next place of the grid: an odd place
% opens a row, an even one closes it.
\newcommand{\GridCurveZone}[2]{%
    \IfFigureDrawnZone{#1}{%
        \stepcounter{gridfigure}%
        \ifodd\value{gridfigure}%
            \ifnum\value{gridfigure}>1 \\[0.5cm]\fi
        \else
            \hfill
        \fi
        \begin{minipage}[t]{0.48\textwidth}
            \centering
            \includegraphics[width=\textwidth]{\Producer_#1_\PCSname.\DTRPcs.\OCname}
            \begin{minipage}[t]{0.8\textwidth}
                \small #2
            \end{minipage}
        \end{minipage}%
    }%
}

\newcommand{\GridCurvesDescriptionZone}{
    \GridCurveDescriptionZone{fig_V}{Voltage at InternalNode1.}
    \GridCurveDescriptionZone{fig_UIt}{Voltage magnitude at InternalNode2 (injector terminal).}
    \GridCurveDescriptionZone{fig_P}{Active power injected by the converter at its control
    point.}
    \GridCurveDescriptionZone{fig_Q}{Reactive power injected by the converter at its control
    point.}
    \GridCurveDescriptionZone{fig_InternalNode1P}{Active power at InternalNode1, where the DTR
    sets the operating point of the test. It is not compared against the reference.}
    \GridCurveDescriptionZone{fig_InternalNode1Q}{Reactive power at InternalNode1, where the DTR
    sets the operating point of the test. It is not compared against the reference.}
    \GridCurveDescriptionZone{fig_Ip}{Active current injected at InternalNode2 (injector
    terminal).}
    \GridCurveDescriptionZone{fig_Iq}{Reactive current injected at InternalNode2 (injector
    terminal).}
    \GridCurveDescriptionZone{fig_I}{Active, reactive, and magnitude currents at InternalNode2
    (injector terminal), with all BESS units displayed on a single plot.}
}

% For now we won't use floats for figures, to get more precise placement
\newcommand{\GridCurvesZone}{%
    \setcounter{gridfigure}{0}%
    \noindent
    \GridCurveZone{fig_V}{Voltage magnitude measured at InternalNode1.}
    \GridCurveZone{fig_UIt}{Voltage magnitude at InternalNode2 (injector terminal).}
    \GridCurveZone{fig_P}{Real power output P, injected by the converter at its control point.}
    \GridCurveZone{fig_Q}{Reactive power output Q, injected by the converter at its control
        point.}
    \GridCurveZone{fig_InternalNode1P}{Active power at InternalNode1, where the DTR sets the
        operating point.}
    \GridCurveZone{fig_InternalNode1Q}{Reactive power at InternalNode1, where the DTR sets the
        operating point.}
    \GridCurveZone{fig_Ip}{Active current output Ip, injected at InternalNode2 (injector
        terminal).}
    \GridCurveZone{fig_Iq}{Reactive current output Iq, injected at InternalNode2 (injector
        terminal).}
    \GridCurveZone{fig_I}{Active (light blue), reactive (violet), and magnitude (yellow)
        currents at InternalNode2 (injector terminal). All BESS units shown on one plot.}
}
%End snippet for the figures of a test